
Comme \(\underline{E}=(\underline{O}_Y)^{r+1}\), on a \(E=r+1\), et par suite :
\[
 \widetilde c(E)=(r+1,1).
\]
Posons :
\[
 x=\cl_{\Gr^1K^{\bullet}(P)}(\xi-1)
 =
 \cl_{\Gr^1K^{\bullet}(P)}(1-\xi^{-1})
 =
 c^1(\xi).
\]
Comme \(\xi^{-1}=\check \xi\), on a :
\[
 \widetilde c(\xi^{-1})=(1,1-x).
\]
On en déduit, d'après V (6.2.1) :
\[
 \widetilde c(T_f)=(r,(1-x)^{r+1}).
\]
Compte tenu de 3.5, on obtient finalement :
\[
 \widetilde c(\check T_f)=(r,(1+x)^{r+1}).
\]

D'après la définition de \(\Todd\), on a donc :
\begin{equation}\tag{5.3.1}
 \Todd(\check T_f)=
 \left(\frac{x}{1-\exp(-x)}\right)^{r+1}.
\end{equation}

\textbf{5.4.} En utilisant la base de \(K^{\bullet}(P)\) sur \(K^{\bullet}(Y)\) formée des éléments \(\xi^k\), avec \(-r\leq k\leq0\), on voit par additivité qu'il suffit de prouver le théorème pour des éléments de la forme \(f^*(a)\cdot\xi^k\). On a alors :
\[
 \ch\bigl(f_*(f^*(a)\xi^k)\bigr)
 =
 \ch(a f_*(\xi^k))
 =
 \ch(a)\cdot\ch(f_*(\xi^k)).
\]
D'autre part :
\[
 f_{\mathrm{gr}}\bigl(\ch(f^*(a)\xi^k)\Todd(\check T_f)\bigr)
 =
 f_{\mathrm{gr}}\bigl(f^{\mathrm{gr}}(\ch(a))\ch(\xi^k)\Todd(\check T_f)\bigr)
 =
 \ch(a)\cdot f_{\mathrm{gr}}\bigl(\ch(\xi^k)\Todd(\check T_f)\bigr).
\]

On voit donc qu'il suffit de montrer le théorème pour les \(\xi^k\), avec \(-r\leq k\leq0\). Or d'après (VI (5.2.3)), on a
\[
 f_*(\xi^k)=0\quad\text{pour }-r\leq k<0~;\qquad f_*(1)=1.
\]
Par suite :
\begin{equation}\tag{5.4.1}
 \ch(f_*(\xi^k))=0\quad\text{si }-r\leq k<0~;\qquad
 \ch(f_*(1))=1.
\end{equation}

